\documentclass[pdflatex,sn-mathphys-num]{sn-jnl}%

\usepackage{graphicx}%
\usepackage{multirow}%
\usepackage{amsmath,amssymb,amsfonts}%
\usepackage{amsthm}%
\usepackage{mathrsfs}%
\usepackage[title]{appendix}%
\usepackage{xcolor}%
\usepackage{textcomp}%
\usepackage{manyfoot}%
\usepackage{booktabs}%
\usepackage{listings}%

\theoremstyle{thmstyleone}%
\newtheorem{theorem}{Theorem}%
\newtheorem{proposition}[theorem]{Proposition}%

\theoremstyle{thmstyletwo}%
\newtheorem{example}{Example}%
\newtheorem{remark}{Remark}%

\theoremstyle{thmstylethree}%
\newtheorem{definition}{Definition}%

\raggedbottom

\usepackage{tabularx}
\usepackage{placeins}
\usepackage{adjustbox}

\newcommand{\manuscriptfigure}[2][]{\includegraphics[#1]{#2}}

\begin{document}

\title[Are 100 Queries Enough?]{Are 100 Queries Enough? Small-Data Adaptation of Open LLMs for Scientific Reranking}

\author[1]{\fnm{Andrei-Răzvan} \sur{Joldea}}\email{andrei.joldea@student.upt.ro}

\author[1]{\fnm{Andrei-Ștefan} \sur{Bulzan}}\email{stefan.bulzan@student.upt.ro}

\author[1]{\fnm{Bogdan} \sur{Morariu}}\email{bogdan.morariu@student.upt.ro}

\author[1]{\fnm{Diana} \sur{Cernăzanu-Glăvan}}\email{diana.cernazanu-glavan@student.upt.ro}

\author*[1]{\fnm{Cosmin} \sur{Cernăzanu-Glăvan}}\email{cosmin.cernazanu@cs.upt.ro}

\affil[1]{\orgdiv{Department of Computer and Information Technology}, \orgname{Universitatea Politehnica Timișoara}, \orgaddress{\city{Timișoara}, \country{Romania}}}

\abstract{
Scientific literature search requires resolving complex queries over large academic
corpora. Large-language-model (LLM) reranking can improve this process, but accurate
systems commonly depend on large proprietary models. This raises concerns about
reproducibility, computational constraints and deployment. We investigate whether
limited supervision is sufficient to adapt open instruction-tuned models, ranging
from 8B to 27B parameters, into effective listwise scientific-paper rerankers.
Starting from 100 human-authored, multi-aspect DORIS-MAE~\cite{Wang} parent queries and their
released relevance labels, we use aspect decomposition, candidate pool subsampling
and pool-level augmentation to expand the dataset to 4,494 training records. In a
nested three-seed learning curve with 10, 25, 50, 75 and 100 parent queries, most of
the Llama 3.1 8B improvement appears by 10--25 parents. The 100-parent condition
obtains the highest Recall@20 (R@20), while Recall@5 (R@5) and normalized
discounted cumulative gain at 20 (nDCG@20) peak at 25 parents. Thus, 100 queries
are sufficient for the observed transfer, but are not shown to be necessary.
Our lexical audit finds no exact or near-duplicate query-text matches with
LitSearch or CSFCube, although some papers are reused across datasets.
Matched controls support aspect-level supervision, but do not establish
an independent benefit from distinct-pool augmentation. A compact supervised
cross-encoder provides a strong practical alternative. Together, these results
support data-efficient adaptation while showing that the gains depend on the
model, metric and evaluation setting.
}

\keywords{scientific literature search, neural reranking, listwise reranking, supervised fine-tuning, large language models, information retrieval}

\maketitle

\section{Introduction}
\label{sec:introduction}
Literature search is a fundamental part of scientific research. It includes
surveying prior work, identifying methods that satisfy specific constraints and
finding relevant papers when only high-level concepts are known. LitSearch
captures this process in a benchmark that pairs realistic research queries with
one or more ground-truth documents. By evaluating both retrieval and downstream
reranking, the benchmark tests whether a model can reason over long scientific
texts~\cite{Ajith}. Its results show that modern embedding-based retrievers
generally outperform sparse retrieval with BM25, while LLM-based reranking
provides further improvements.

However, accurate reranking pipelines often rely on large proprietary models,
such as GPT-class rerankers, and on zero-shot prompting. Their effectiveness can
vary with the prompt design, candidate-list length and positional bias~\cite{Sun}.
Our primary objective is to balance two requirements: the model should be small
enough to self-host without external APIs, but still competitive for scientific
literature reranking. We therefore select open instruction-tuned models from
the Gemma 3~\cite{gemma3}, Qwen2.5~\cite{qwen25} and Llama 3~\cite{llama3}
families, whose long contexts support prompts containing several candidate abstracts.

We organize this work around one research question: can a small but structured
scientific retrieval dataset be used to adapt open instruction-tuned LLMs into
useful listwise rerankers that transfer to different scientific search benchmarks?
Our contributions are related to methodology and system design:
\begin{itemize}
\item We combine aspect decomposition, candidate subsampling and set-unique pool
augmentation to turn DORIS-MAE's 100 human-authored parent queries into 4,494
training lists. We train on its released relevance labels, which were produced
using ChatGPT~\cite{Wang}, without generating additional teacher rankings.
\item We directly evaluate sample efficiency through nested parent-query budgets
and three training seeds. Matched ablations, paired uncertainty estimates and
document-overlap controls qualify the contribution of the data-construction recipe.
\item We evaluate transfer to LitSearch and CSFCube, compare with open reranking
baselines and an external GPT-4o reference, and assess a compact supervised
cross-encoder's ranking quality and measured inference cost.
\end{itemize}
LitSearch uses natural-language literature questions, whereas CSFCube follows a
faceted query-by-example paradigm. Evaluation on both reduces dependence on a
single query format. Figure~\ref{fig:pipeline} summarizes the workflow. 
Online
Resource~1 provides the detailed protocols and supporting experiments.

\section{Related Work}

Scientific paper retrieval has led to datasets and benchmarks with different query
forms and evaluation protocols. DORIS-MAE focuses on complex, human-authored,
multi-aspect queries in computer science. For each query, it provides a candidate
pool with graded relevance on a 0-2 scale \cite{Wang}. LitSearch instead studies
literature search questions and pairs every query with one or more ground-truth
papers. Some queries are generated by an LLM from inline citations, while others
are written by researchers about their own work \cite{Ajith}. CSFCube uses a third
format: faceted query-by-example. A query consists of a scientific paper and a
specific facet, such as background, method or result, and relevance depends on
similarity along that facet. The benchmark uses depth-$k$ pooling under
TREC-style practices to obtain graded relevance judgments \cite{Mysore}.

Most retrieval pipelines contain two stages: an initial retrieval step followed by
reranking. Sparse lexical methods such as BM25 remain strong and efficient
reference points. Neural retrieval, however, has developed in several directions.
Dual encoders such as Dense Passage Retrieval (DPR) \cite{Karpukhin} support
efficient dense retrieval, while late-interaction models such as ColBERT
\cite{Khattab} exchange larger indexes for finer-grained query-document matching.
SciBERT \cite{Beltagy} uses language-model pretraining on scientific
text, while SPECTER \cite{Cohan} learns scientific document
representations using citation-based supervision. Broad evaluations
such as Benchmarking-IR (BEIR) \cite{Thakur} and the Massive Text Embedding
Benchmark (MTEB) \cite{Muennighoff} show why cross-task and cross-domain testing
is necessary: strong in-domain performance does not always transfer to a new
setting.

Reranking methods are commonly divided into pointwise, pairwise and listwise
approaches, according to whether the loss is defined over individual documents,
document pairs or complete ranked lists. Listwise methods such as ListMLE
\cite{Xia} optimize a probability over permutations and form the basis of much of the
listwise reranking literature. For neural reranking, monoT5 \cite{Nogueira}
reformulated the task as token generation: T5 is fine-tuned to emit ``true'' or
``false,'' and relevance is scored by the probability of ``true'' after softmax
normalization over these two token logits. RankT5
\cite{Zhuang} added pairwise and listwise losses to T5-based reranking, while RankGPT \cite{Sun} adopted a generative formulation
toward full listwise ranking. These systems are usually evaluated with top-weighted
IR metrics that support graded relevance, such as (n)DCG.

More recent work uses LLMs directly as rerankers. In RankGPT \cite{Sun}, for
example, the model returns an ordering of candidate identifiers for a given query.
This line of research has increased interest in open-weight listwise
rerankers. RankZephyr \cite{Pradeep} aims to approach proprietary-model
performance, while enabling inference without relying on proprietary APIs. Other work
addresses efficiency by reducing decoding or constraining the ranking interface.
FIRST, for instance, performs listwise reranking through single-token decoding and
a learning-to-rank objective \cite{Reddy}. Surveys of LLMs for IR discuss
deployment costs and model adaptation \cite{Zhu}. These concerns motivate our
study of reranking with limited task-specific supervision.

\section{Methodology}
\label{sec:methodology}

We study listwise scientific paper reranking. Given a query $q$ and an initial set
of $N$ candidate documents $\{d_1, \dots, d_N\}$, we ask the model to predict a
permutation that orders them by decreasing relevance. Each document contains a
paper title and abstract. This formulation follows earlier LLM-based listwise
prompts and benchmark implementations, including RankGPT-style \cite{Sun}
prompts that request an ordered sequence such as ``[2] $>$ [5] $>$ \dots''.

\begin{figure}[t]
\centering
\manuscriptfigure[width=\textwidth]{architecture.drawio.pdf}
\caption{Our scientific paper reranking pipeline: DORIS-MAE data construction and fine-tuning, followed by evaluation on LitSearch and CSFCube}
\label{fig:pipeline}
\end{figure}

\subsection{Data preparation}

We construct our training set from DORIS-MAE \cite{Wang}. The released dataset
contains 100 complex, human-authored queries from several computer science
subfields. Each query is divided into aspects and, when available, sub-aspects. It
is also paired with a pool of approximately 99–138 paper abstracts and graded
relevance labels. The underlying corpus contains 363,133 computer science paper
abstracts from arXiv. We use this corpus both to
construct candidate pools and to retrieve paper text when formatting prompts.

We do not treat a complete multi-aspect query as one atomic training unit.
Instead, every aspect or sub-aspect sentence becomes an ``aspect-level query'' and
retains its parent's candidate pool and relevance labels. We perform the extraction
as follows:

\begin{itemize}
\item \textbf{Aspect-text:} We use the provided natural-language sentences for
main aspects and the short descriptors for sub-aspects.
\item \textbf{Relevance labels:} We extract each label from the per-(abstract,
aspect) score table when it is available. If a given (abstract, aspect) pair has no
annotation, we assign it a score of 0.
\end{itemize}

This operation transforms the 100 original queries into 1,501 aspect-level ranking
instances. We then remove instances for which the entire pool has a score of zero,
leaving 1,498 valid records.

\subsection{Candidate pool subsampling}

Managing context length is a practical challenge in listwise LLM reranking. A
prompt containing 100 abstracts can produce a training sequence that exceeds the
model's context window. Even when it fits, transformer attention and KV-cache
growth can cause out-of-memory errors. Sliding windows address input-length
constraints by reranking subsets of candidates \cite{Sun}. Single-token decoding
instead reduces output-generation cost \cite{Reddy}.

To fine-tune efficiently on one GPU, we sample training lists with a target size
$K$, denoted by \texttt{TARGET\_SIZE}, from each aspect-level candidate set.
The parameter \texttt{MAX\_POSITIVES} limits the initial positive selection.
In our experiments, $K = 50$ and \texttt{MAX\_POSITIVES} $= 5$. Negatives with
score 0 then fill the remaining positions. These are sampling parameters rather
than strict final constraints: a fallback positive can increase the list size to
51, while replenishment from the remaining pool can exceed the positive limit.
Reconstructing the standard recipe gives 4,312 lists of length 50 and 182 of
length 51. 1,392 lists contain more than five positives. We retain this historical
behavior in the reported experiments. We then shuffle the candidates to produce a
realistic input order. For larger models such as Gemma 3 27B, we further control
the prompt length by truncating every candidate passage to a bounded prefix of
approximately 300 words.

This procedure resembles standard training with hard negatives, but there is an
important difference: we draw negatives directly from the existing candidate pool
instead of mining them explicitly for difficulty. We construct the target output by
sorting candidates in descending order of relevance. When scores are tied, the
candidate's index in the presented pool determines the order. Identical pools
therefore produce identical targets across training epochs.

We also audit whether missing judgments were implicitly converted to relevance
zero during dataset construction. Against the original DORIS-MAE annotations,
all 163,969 candidate--aspect pairs in the 1,498 retained training rows have an
explicit judgment, including all 88,285 score-zero pairs. The missing-to-zero
fallback was therefore never triggered in the reported dataset. This audit rules
out unjudged candidates being treated as negatives, but does not remove possible
sensitivity to deterministic ordering within tied relevance grades.
We test this sensitivity separately in Section~\ref{sec:tie-sensitivity}, using
matched deterministic and epoch-randomized targets.

\subsection{Pool-level augmentation with set-level uniqueness}

We vary candidate composition by exposing the model to different candidate
combinations. We generate up to $V$ distinct pools for each aspect-level
query, with $V = 3$ in our experiments. Candidates are divided into strong
positives (score 2), weak positives (score 1) and negatives (score 0). During
sampling, we cycle through four positive-selection modes:
\begin{itemize}
\item \textit{Only strong positives:} We draw positives exclusively from score-2
documents.
\item \textit{Only weak positives:} We draw positives exclusively from score-1
documents.
\item \textit{Mixed:} We combine score-1 and score-2 documents.
\item \textit{Any positives:} We sample from the complete positive set, regardless
of grade.
\end{itemize}
We use two safeguards during generation. First, we enforce \emph{set-level
uniqueness}. Two pools with the same document identifiers are treated as
duplicates, irrespective of internal order, and we retain only the first. Thus, the
augmentation changes candidate composition rather than merely permuting
positions. Second, we limit the number of rejected sampling attempts for each
query. A query can consequently produce fewer than $V$ variants when its
candidate pool is too small to yield enough unique sets. Finally, we shuffle every
selected list before prompt formatting to reduce positional bias. Together, these
steps expand the 1,498 filtered aspect-level queries into 4,494 distinct training
instances.

\subsection{Parent-query supervision-budget protocol}
\label{sec:parent-budget-protocol}

We measure sample efficiency at the level of independent human-authored parent
queries rather than treating the correlated aspect rows or augmented lists as
independent supervision. We use Llama 3.1 8B as the representative model and
train budgets $B\in\{10,25,50,75,100\}$ with seeds 42, 2026 and 3407. For each
seed, we deterministically shuffle parents within each of the five DORIS-MAE
query categories and take $B/5$ parents from every category. Smaller budgets are
prefixes of the same category-specific orders, so the subsets are nested within a
seed. Selection occurs before aspect decomposition: all eligible aspect rows from
a selected parent are retained together, so a parent's aspect rows are never
split across inclusion groups within a run.

Because parents contain different numbers of eligible aspects, the selected
aspect counts vary across seeds. For budgets 10, 25, 50, 75 and 100, the three
seed-specific counts are respectively $(144,147,160)$, $(348,406,389)$,
$(705,782,753)$, $(1108,1137,1124)$ and $(1498,1498,1498)$. We apply the
standard recipe unchanged: at most three set-unique pools per aspect,
$K_{\mathrm{train}}=50$, $P=5$, two epochs, batch size two and no gradient
accumulation. This produces 432--480, 1044--1218, 2115--2346, 3324--3411 and
4494 augmented lists across the five budgets. The corresponding nominal
optimizer-step ranges are 432--480, 1044--1218, 2116--2346, 3324--3412 and
4494. We evaluate all 15 checkpoints on the same 597 LitSearch queries at
$K_{\mathrm{eval}}=100$ with BM25, GRIT and Qwen3 retrieval. Within each seed
we first average each metric over the three retrievers, then report the mean and
sample standard deviation across seeds. Retriever-specific results are retained
in Online Resource~1, Table~S11.

\subsection{Leakage and overlap analysis}

All three datasets originate from scientific-paper corpora, so we explicitly
measure overlap between the DORIS-MAE-derived fine-tuning data and the
LitSearch / CSFCube evaluation material. We examine query-level and
document-level overlap separately. To match documents, we compare normalized
titles, normalized title-abstract strings, available arXiv and DOI identifiers and
dataset-specific source identifiers. We do not treat raw integer IDs as conclusive
across datasets because each dataset uses a different ID namespace.

We find no exact or near-duplicate query overlap between the DORIS-MAE training
queries and the LitSearch or CSFCube queries. Document-level overlap does exist.
As Table~\ref{tab:leakage-overlap} shows, DORIS-MAE prompt construction exposes
8,480 unique documents. Of these, 123 have exact normalized title-abstract
matches in the CSFCube evaluation documents, while 227 match documents in the
LitSearch corpus. When we restrict the comparison to positive DORIS-MAE prompt
candidates, the corresponding counts are 108 for CSFCube and 202 for LitSearch.
We therefore interpret the primary cohort as evidence of query-level
generalization across benchmarks. Section~\ref{sec:document-disjoint} supplements
this analysis with a strict LitSearch rerun.

\begin{table}[htbp]
\centering
\footnotesize
\setlength{\tabcolsep}{3pt}
\begin{adjustbox}{max width=\linewidth}
\begin{tabular}{p{0.46\columnwidth} r r r}
\hline
Comparison & Left & Right & Matches \\
\hline
DORIS prompt docs vs CSFCube eval & 8,480 & 4,205 & 123 \\
DORIS positive prompt docs vs CSFCube eval & 7,216 & 4,205 & 108 \\
DORIS prompt docs vs LitSearch corpus & 8,480 & 64,183 & 227 \\
DORIS positive prompt docs vs LitSearch corpus & 7,216 & 64,183 & 202 \\
\hline
\end{tabular}
\end{adjustbox}
\caption{Exact normalized title-abstract matches between DORIS-MAE exposure and evaluation corpora. No exact or near-duplicate query overlap was found}
\label{tab:leakage-overlap}
\end{table}

\section{Experimental Setup}
\label{sec:setup}
\subsection{Prompting, training and output repair}
Each training example is a chat-style conversation. The user asks the model to
order indexed title-abstract passages by relevance. The assistant response is
the target permutation. We compute loss only on assistant-response tokens.
The complete prompt is given in Online Resource~1. At LitSearch inference, all
100 candidates are placed in one prompt, indexed from $[1]$ to $[100]$, with
each passage truncated to approximately 300 words. The main generative rerankers
do not use sliding windows. However, the reranking inference code includes a safety mechanism which progressively reduces the number of included documents in the prompt in case of failures.

For the primary cohort, Llama 3.1 8B, Gemma 3 12B and Qwen 2.5 14B use 16-bit
LoRA~\cite{Hu}; Gemma 3 27B uses 4-bit QLoRA~\cite{Dettmers}. We apply LoRA to
the attention and feed-forward linear projections, with rank $r=128$,
$\alpha=256$ and dropout $0.05$. Vision encoders remain frozen. Training uses
\texttt{bfloat16}, Unsloth~\cite{unsloth} and TRL's \texttt{SFTTrainer}, with
AdamW, a cosine schedule, 20 warmup steps, weight decay $0.01$, learning rate
$1\times10^{-5}$ and two epochs. Context length is capped at 16,384 tokens.
Table~\ref{tab:model-training-settings} gives model-specific settings.
Adapters are merged into 16-bit checkpoints for local inference.

\begin{table}[t]
\centering
\scriptsize
\setlength{\tabcolsep}{4pt}
\renewcommand{\arraystretch}{1.08}
\begin{tabularx}{\textwidth}{p{0.16\textwidth} X c c c c}
\hline
Model & Base checkpoint and loading mode & Batch & Accum. & Effective batch & Epochs \\
\hline
Llama 3.1 8B & \texttt{meta-llama/\allowbreak Llama-3.1-8B-Instruct}; 16-bit LoRA & 2 & 1 & 2 & 2 \\
Gemma 3 12B & \texttt{google/gemma-3-12b-it}; 16-bit LoRA & 2 & 1 & 2 & 2 \\
Qwen 2.5 14B & \texttt{Qwen/Qwen2.5-14B-Instruct}; 16-bit LoRA & 2 & 1 & 2 & 2 \\
Gemma 3 27B & \texttt{google/gemma-3-27b-it}; 4-bit QLoRA & 2 & 1 & 2 & 2 \\
\hline
\end{tabularx}
\caption{Primary-cohort fine-tuning settings. Batch and Accum. denote per-device batch size and gradient accumulation; all runs use one GPU}
\label{tab:model-training-settings}
\end{table}

We extract bracketed identifiers, discard out-of-range IDs and retain only the
first occurrence of each candidate. Missing candidates are appended in their
original first-stage order; if no usable IDs are produced, we retain the complete
original order. Metrics therefore describe the model together with this repair
rule. We examine raw validity separately in the tie-policy experiment and
compute benchmark.

\subsection{Compact supervised reranker}
\label{sec:compact-training}
We also adapt GTE-ModernBERT~\cite{gte-modernbert}, a cross-encoder with
149,605,633 parameters initialized from the ranking-pretrained checkpoint
\texttt{Alibaba-NLP/gte-reranker-modernbert-base}, revision \texttt{204bdd4}.
Here, ``zero-shot'' means without additional DORIS-MAE adaptation, not without
prior ranking supervision. Both released and adapted versions are evaluated.

The compact model uses 4,494 frozen DORIS-MAE-derived lists with the same nominal
sampling recipe. Its artifact exposes 8,477 unique document IDs, compared with
8,480 in the reconstructed LLM preparation data. Exact identity with every
historical LLM training list is therefore not established, and this is not a
fully matched training-data ablation. We update all parameters in
\texttt{bfloat16}, using LambdaLoss with NDCGLoss2++ weighting at $k=20$, seed
3407, two epochs and learning rate $2\times10^{-5}$. Training passages retain
the 300-word prefix and query-document pairs are limited to 1,024 tokens.
Hyperparameters are fixed in advance; neither evaluation benchmark is used for
checkpoint selection. Complete settings are retained in Online Resource~1.

At evaluation, GTE scores and sorts candidates. On LitSearch it receives the
same frozen top-100 lists as the LLMs. On CSFCube, it scores every candidate in
the released pool, representing the query paper by its title and full abstract
subject to the pair limit.

\subsection{Benchmarks and evaluation metrics}
We extend the released LitSearch harness with local vLLM inference~\cite{Kwon}
and Qwen3 embeddings~\cite{Zhang}. LitSearch~\cite{Ajith} contains 597 queries
and 64,183 ML and NLP papers. We retrieve 200 candidates and rerank the top 100
using BM25, GritLM-7B (GRIT)~\cite{gritlm} and Qwen3-Embedding-8B (Qwen3).
CSFCube~\cite{Mysore} contains 50 query-facet cases from 34 unique query papers,
with released candidate pools and graded relevance for background, method and
result. It evaluates reranking alone.

LitSearch uses binary relevance. We compute
\(R@k=|\text{top-}k\text{ retrieved}\cap\text{relevant}|/|\text{relevant}|\)
and macro-average across queries. For nDCG, relevant documents receive gain 1
and others gain 0, with logarithmic discount $1/\log_2(i+1)$ and normalization
by the ideal ranking. R@20 and nDCG@20 are primary, while R@5 is a descriptive
secondary endpoint. Our recall implementation reproduces published retrieval-only
results to within rounding (Online Resource~1, Table~S2).

CSFCube uses grades 0--3 directly as gains. Its nDCG\%@20 evaluates the top
20\% of each candidate pool, at cutoff $\lfloor(20/100)n\rfloor$ for pool size
$n$, rather than a fixed rank. We retain the released evaluator. Its dev aggregate
contains 24 cases; test is the equal-weight mean of the 26- and 24-case test
folds. Dev cases reappear in one test fold, so these aggregates are descriptive,
not independent evaluations. Paired inference instead pools all 50 cases into a
case-weighted mean and clusters by the 34 unique papers.

We use 10,000 paired bootstrap replicates for pointwise 95\% percentile intervals
and 10,000 two-sided sign-flip permutations with the plus-one correction.
LitSearch resamples queries. CSFCube resamples papers with their facets kept
together. Holm correction is applied within each metric and prespecified
comparison family. The 12 LitSearch base--FT, six local-baseline, two
pool-augmentation and four CSFCube base--FT contrasts form separate families.
Compact-model comparisons use separate families detailed in Online Resource~1.
Tie-policy families are described in Section~\ref{sec:tie-sensitivity}.
Intervals are not multiplicity-adjusted and these tests
measure evaluation uncertainty conditional on retained checkpoints. Training-seed
variation is reported separately.

\subsection{Experiment provenance and model selection}
\label{sec:reproducibility}
The primary and reconstructed LLM experiments use one NVIDIA H200 NVL. We keep
these cohorts separate because software changes prevented exact reconstruction
of the original environment. Ablations compare against a newly trained
standard-recipe control in the same reconstructed environment; all reconstructed
models use 16-bit LoRA, including Gemma 3 27B. We do not interpret differences
between environments as treatment effects. Online Resource~1 records software
versions, seeds, manifests and a complete four-model replication.

The historical hyperparameter search used an 80/10/10 split at formatted-list
level, so aspects and augmented pools from one parent could cross partitions.
We therefore perform an additional Llama 3.1 8B learning-rate study with
80/10/10 parent-query-disjoint partitions, comparing $5\times10^{-6}$,
$1\times10^{-5}$ and $2\times10^{-5}$ across three split assignments. Selection
minimizes response-token negative log-likelihood averaged equally over validation
parents and then splits. Internal-test losses are evaluated after selection, and
the chosen configuration is refitted on all 4,494 lists. This study uses a
separate environment and does not replace historical checkpoints. LitSearch
and CSFCube remain external evaluation datasets throughout.

\section{Results}
\label{sec:results}
``Base'' denotes an off-the-shelf instruction-tuned model evaluated zero-shot
with our ranking prompt; ``FT'' denotes its adapted checkpoint. Within each
local base--FT pair, the first-stage candidate pool and input order are fixed.
We report the primary cohort first and identify reconstructed or separately
trained controls explicitly.
We note that all differences for the locally evaluated results are calculated from unrounded
scores and rounded only for presentation. Therefore, they may differ slightly
from differences between displayed values.

\subsection{LitSearch transfer}
Table~\ref{tab:litsearch-main} presents the primary results and
Table~\ref{tab:litsearch-bootstrap} their paired uncertainty. Without reranking,
Qwen3 obtains overall R@20 of $80.5\%$, R@5 of $69.0\%$ and nDCG@20 of
$62.9\%$. Fine-tuning generally improves the mid-size models. The largest
R@5 change appears for Llama 3.1 8B with Qwen3, from $52.2\%$ to $68.4\%$
($+16.2$ pp), approaching the substantially larger base Gemma 3 27B model.

\begin{table}[ht]
\centering
\renewcommand{\arraystretch}{1.1}
\setlength{\tabcolsep}{3.5pt}
\footnotesize
\begin{adjustbox}{max width=\linewidth}
\begin{tabular}{l l c c c c c c}
\hline
\multirow{2}{*}{Model} & \multirow{2}{*}{Retriever}
 & \multicolumn{2}{c}{R@20} & \multicolumn{2}{c}{R@5}
 & \multicolumn{2}{c}{nDCG@20} \\
 & & base & FT & base & FT & base & FT \\
\hline
\multirow{3}{*}{Llama 3.1 8B}
 & BM25  & 54.6 & \textbf{66.2} & 39.0 & \textbf{51.4} & 38.7 & \textbf{47.8} \\
 & GRIT  & 72.1 & \textbf{79.4} & 54.3 & \textbf{67.6} & 52.4 & \textbf{62.6} \\
 & Qwen3 & 72.3 & \textbf{79.5} & 52.2 & \textbf{68.4} & 50.7 & \textbf{63.4} \\
\hline
\multirow{3}{*}{Gemma 3 12B}
 & BM25  & 60.0 & \textbf{66.1} & 51.4 & \textbf{53.0} & 47.5 & \textbf{47.8} \\
 & GRIT  & 79.9 & \textbf{81.2} & 69.7 & \textbf{70.2} & 63.1 & \textbf{64.1} \\
 & Qwen3 & 77.6 & \textbf{81.2} & 69.4 & \textbf{71.0} & 61.1 & \textbf{64.6} \\
\hline
\multirow{3}{*}{Qwen 2.5 14B}
 & BM25  & 60.6 & \textbf{65.1} & \textbf{51.8} & 51.3 & 46.8 & \textbf{47.7} \\
 & GRIT  & 77.6 & \textbf{80.7} & 70.3 & \textbf{71.0} & 61.0 & \textbf{64.1} \\
 & Qwen3 & 78.7 & \textbf{80.4} & 69.0 & \textbf{70.5} & 61.9 & \textbf{64.2} \\
\hline
\multirow{3}{*}{Gemma 3 27B (4-bit)}
 & BM25  & 65.8 & \textbf{66.2} & \textbf{56.4} & 56.3 & \textbf{51.9} & 49.7 \\
 & GRIT  & 81.7 & \textbf{82.8} & \textbf{71.8} & 70.9 & \textbf{65.7} & 65.0 \\
 & Qwen3 & 79.5 & \textbf{84.0} & 70.9 & \textbf{72.5} & 65.2 & \textbf{65.9} \\
\hline
\end{tabular}
\end{adjustbox}
\caption{Primary-cohort LitSearch scores over 597 queries at $K_{\mathrm{eval}}=100$. Base and FT use the same first-stage candidates; FT denotes listwise fine-tuning}
\label{tab:litsearch-main}
\end{table}
\begin{table}[t]
\centering
\scriptsize
\setlength{\tabcolsep}{4pt}
\renewcommand{\arraystretch}{0.96}
\begin{adjustbox}{max width=\linewidth}
\begin{tabular}{l l l r c c c}
\hline
Model & Retriever & Metric & $\Delta$ & 95\% CI & $p$ & $p_{\mathrm{H}}$ \\
\hline
Llama 3.1 8B & BM25 & R@20 & 11.67 & [8.57, 14.91] & $<.001$ & .0012 \\
Llama 3.1 8B & BM25 & nDCG@20 & 9.08 & [6.74, 11.47] & $<.001$ & .0012 \\
Llama 3.1 8B & GRIT & R@20 & 7.29 & [4.19, 10.30] & $<.001$ & .0012 \\
Llama 3.1 8B & GRIT & nDCG@20 & 10.20 & [7.38, 13.00] & $<.001$ & .0012 \\
Llama 3.1 8B & Qwen3 & R@20 & 7.22 & [4.40, 10.05] & $<.001$ & .0012 \\
Llama 3.1 8B & Qwen3 & nDCG@20 & 12.72 & [10.06, 15.40] & $<.001$ & .0012 \\
\hline
Gemma 3 12B & BM25 & R@20 & 6.11 & [3.69, 8.54] & $<.001$ & .0012 \\
Gemma 3 12B & BM25 & nDCG@20 & 0.33 & [$-$1.42, 2.09] & .7144 & 1.0000 \\
Gemma 3 12B & GRIT & R@20 & 1.31 & [$-$1.03, 3.71] & .3014 & .9041 \\
Gemma 3 12B & GRIT & nDCG@20 & 0.97 & [$-$1.02, 2.98] & .3435 & 1.0000 \\
Gemma 3 12B & Qwen3 & R@20 & 3.57 & [1.23, 5.89] & .0035 & .0210 \\
Gemma 3 12B & Qwen3 & nDCG@20 & 3.49 & [1.55, 5.43] & $<.001$ & .0027 \\
\hline
Qwen 2.5 14B & BM25 & R@20 & 4.52 & [2.12, 6.95] & $<.001$ & .0049 \\
Qwen 2.5 14B & BM25 & nDCG@20 & 0.87 & [$-$0.90, 2.68] & .3454 & 1.0000 \\
Qwen 2.5 14B & GRIT & R@20 & 3.11 & [0.84, 5.53] & .0091 & .0455 \\
Qwen 2.5 14B & GRIT & nDCG@20 & 3.09 & [0.86, 5.38] & .0062 & .0496 \\
Qwen 2.5 14B & Qwen3 & R@20 & 1.79 & [$-$0.36, 3.85] & .1080 & .4320 \\
Qwen 2.5 14B & Qwen3 & nDCG@20 & 2.30 & [0.38, 4.24] & .0197 & .1182 \\
\hline
Gemma 3 27B & BM25 & R@20 & 0.42 & [$-$1.73, 2.54] & .7212 & .9041 \\
Gemma 3 27B & BM25 & nDCG@20 & $-$2.20 & [$-$3.92, $-$0.47] & .0129 & .0903 \\
Gemma 3 27B & GRIT & R@20 & 1.10 & [$-$1.13, 3.43] & .3676 & .9041 \\
Gemma 3 27B & GRIT & nDCG@20 & $-$0.67 & [$-$2.35, 0.97] & .4444 & 1.0000 \\
Gemma 3 27B & Qwen3 & R@20 & 4.57 & [2.28, 6.85] & $<.001$ & .0016 \\
Gemma 3 27B & Qwen3 & nDCG@20 & 0.67 & [$-$1.08, 2.44] & .4592 & 1.0000 \\
\hline
\end{tabular}
\end{adjustbox}
\caption{LitSearch FT-minus-base differences in percentage points, with pointwise paired 95\% bootstrap intervals. $p$ and $p_{\mathrm{H}}$ denote raw and Holm-adjusted sign-flip values; R@5 intervals are in Online Resource~1}
\label{tab:litsearch-bootstrap}
\end{table}

The type of gain varies with retriever strength. R@20 improves most under BM25, especially for the mid-sized models,
where the reranker can filter more low-quality candidates. R@5 and nDCG@20
benefit most with GRIT and Qwen3, where ordering relevant candidates precisely
becomes more important. All six primary Llama gains remain significant after
Holm correction. For Gemma 3 12B, both primary metrics improve significantly
with Qwen3, while only R@20 does so with BM25. For Qwen 2.5 14B, both GRIT
gains and the BM25 R@20 gain remain significant.

Gemma 3 27B behaves differently: R@20 increases with every retriever, but early
ranking metrics change marginally or decline with BM25 and GRIT. Only its Qwen3
R@20 gain remains significant after correction. Its BM25 nDCG@20 interval
excludes zero before correction, but $p_{\mathrm{H}}=.0903$. We therefore
characterize the larger model as competitive and configuration-dependent.

With Qwen3 candidates, base Gemma 3 27B reaches R@20 of $79.5\%$, slightly
below Qwen3 retrieval alone ($80.5\%$). Fine-tuning raises it to $84.0\%$:
$+3.5$ pp over retrieval and $+4.6$ pp over its zero-shot behavior. A strong
instruction-tuned model can therefore degrade a strong retriever before
adaptation. Detailed subset results and the gain visualization are retained in
Online Resource~1.

\subsection{Supervision-budget learning curve}
\label{sec:parent-learning-curve}
Figure~\ref{fig:parent-query-learning-curve} makes the independent supervision
budget a primary experimental variable. The panels show the same 15 runs against
parent queries, eligible aspects and nominal optimizer steps. The sharpest change
occurs between the zero-shot model and the 10-parent condition. Beyond 10--25
parents, the curves are approximately flat and are not monotonic.

\begin{figure}[ht]
  \centering
  \manuscriptfigure[width=\textwidth]{fig_parent_query_learning_curve.pdf}
  \caption{Llama 3.1 8B learning curve at $K_{\mathrm{eval}}=100$. Filled markers and bars show the mean and sample standard deviation over three seeds after retriever averaging; faint markers show individual runs. Open markers are the zero-shot reference. All panels represent the same runs}
  \label{fig:parent-query-learning-curve}
\end{figure}

The 100-parent model has the highest mean R@20, but improves it by only 1.03
percentage points over 10 parents and 0.44 points over 25. The 25-parent condition
has the highest R@5 and nDCG@20, exceeding the 100-parent condition by 1.03 and
0.70 points, respectively. Increasing aspect examples and optimization steps by
approximately a factor of ten does not produce a corresponding performance
increase. Exact values and retriever-specific results are in
Online Resource~1, Table~S4 and Online Resource~1, Table~S11.

These results answer the title question narrowly. One hundred DORIS-MAE parent
queries are sufficient for the reported transfer, but the curve does not show
that all 100 are necessary. Most measured gain appears with 10--25 parents.
Three seeds do not support a high-power equivalence test, and the best budget
depends on whether the objective is deeper recall or early ranking quality.

\subsection{Parent-grouped hyperparameter selection}
\label{sec:grouped-hpo}
After performing a parent-grouped search, we select $5\times10^{-6}$, which gives the lowest
parent-averaged validation NLL on all three splits. Only the learning rate is
varied. We refit the base model on all 4,494 lists for two epochs at this rate. The resulting checkpoint is separate from the primary cohort. Validation NLLs,
split variation and the shared-parent internal-test results are reported in
Online Resource~1, Section~S14.

\begin{table}[t]
\centering
\small
\setlength{\tabcolsep}{4pt}
\begin{adjustbox}{max width=\linewidth}
\begin{tabular}{lrrrr}
\hline
Retriever & R@20 & R@5 & nDCG@20 & nDCG@5 \\
\hline
BM25  & 66.31 & 55.11 & 50.03 & 46.66 \\
GRIT  & 80.88 & 69.13 & 64.23 & 60.71 \\
Qwen3 & 81.09 & 69.21 & 64.40 & 60.84 \\
\hline
\end{tabular}
\end{adjustbox}
\caption{LitSearch scores for the separately trained, grouped-selected Llama 3.1 8B refit: percentages over 597 queries at $K_{\mathrm{eval}}=100$}
\label{tab:grouped-hpo-litsearch}
\end{table}

Table~\ref{tab:grouped-hpo-litsearch} reports the refit's LitSearch scores.
Its official aggregate CSFCube test scores are R@$20=40.83\%$,
nDCG@$20=43.51\%$ and nDCG\%@$20=43.20\%$.

\subsection{Configuration and matched ablations}
\label{sec:revision-ablations}
The Llama 3.1 8B ablations use the reconstructed environment and a newly trained
standard-recipe control. Table~\ref{tab:litsearch-ablation-main} averages each
metric across the same 597 queries and three retrievers at
$K_{\mathrm{eval}}=100$. Removing aspect decomposition causes the largest
aggregate decline. However, it changes both the number of examples and the
construction of relevance labels, so it measures the complete transformation
rather than query wording alone. Of the tested list sizes and positive caps,
$K_{\mathrm{train}}=50$ and $P=5$ give the best aggregate balance.

\begin{table}[ht]
\centering
\scriptsize
\setlength{\tabcolsep}{3.2pt}
\renewcommand{\arraystretch}{1.08}
\begin{adjustbox}{max width=\linewidth}
\begin{tabular}{l c c c}
\hline
Configuration & R@20 & R@5 & nDCG@20 \\
\hline
Standard: $V=3,K_{\mathrm{train}}=50,P=5$ & \textbf{77.34} & 64.93 & \textbf{60.21} \\
$V=1,K_{\mathrm{train}}=50,P=5$ & 76.39 & \textbf{65.72} & 60.08 \\
$V=3,K_{\mathrm{train}}=20,P=5$ & 75.99 & 65.65 & 59.84 \\
$V=3,K_{\mathrm{train}}=100,P=5$ & 73.41 & 61.56 & 57.98 \\
$V=3,K_{\mathrm{train}}=50,P=1$ & 75.90 & 64.05 & 59.16 \\
$V=3,K_{\mathrm{train}}=50,P=10$ & 75.34 & 63.94 & 59.05 \\
No pool augmentation & 74.99 & 63.95 & 60.01 \\
No aspect decomposition & 72.82 & 60.65 & 55.33 \\
\hline
\end{tabular}
\end{adjustbox}
\caption{Llama 3.1 8B configuration ablations at $K_{\mathrm{eval}}=100$. Scores are percentages averaged over BM25, GRIT and Qwen3; bold marks the column maximum}
\label{tab:litsearch-ablation-main}
\end{table}

The three-pool recipe improves R@20 relative to the strongest $V=1$ control by
$+0.95$ pp, with 95\% CI $[-0.06,1.95]$ and $p_{\mathrm{H}}=.0717$.
This is not significant after correction. Its gain over no pool augmentation
is $+2.35$ pp, with CI $[1.29,3.40]$ and $p_{\mathrm{H}}<.001$.
The no-augmentation condition draws one \texttt{any}-positive pool without
cycling through structured modes. Because these configurations change training
examples and updates, we do not treat the differences as a fully matched causal
decomposition.

We therefore add example-count-matched controls at seed 42. Every condition
contains 4,494 examples and 4,494 nominal optimizer steps. The no-augmentation
condition repeats one pool three times per aspect. The no-aspect condition
repeats its 300 full-query lists 14--15 times. Table~\ref{tab:matched-main-ablation}
does not support an independent benefit from distinct-pool augmentation:
repeating one pool is slightly better on all three aggregate metrics. Removing
aspect decomposition still reduces R@20, R@5 and nDCG@20 by 3.10, 3.70 and
3.87 points. This strengthens the evidence for aspect-level supervision. The
control remains single-seed and does not exactly match non-padding tokens, so
small augmentation differences are descriptive.

\begin{table}[ht]
\centering
\footnotesize
\setlength{\tabcolsep}{3.2pt}
\renewcommand{\arraystretch}{1.08}
\begin{adjustbox}{max width=\linewidth}
\begin{tabular}{l c c c}
\hline
Matched condition & R@20 & R@5 & nDCG@20 \\
\hline
Full recipe & 76.03 & 64.58 & 59.41 \\
No pool augmentation & \textbf{76.88} & \textbf{65.23} & \textbf{59.69} \\
No aspect decomposition & 72.93 & 60.88 & 55.54 \\
\hline
\end{tabular}
\end{adjustbox}
\caption{Example-count-matched Llama 3.1 8B controls at seed 42: 4,494 training examples per condition. Scores are percentages averaged over the three retrievers}
\label{tab:matched-main-ablation}

\end{table}

\subsection{Tie-handling sensitivity}
\label{sec:tie-sensitivity}
We compare deterministic input-order tie breaking with epoch-wise randomization
within relevance grades. Both policies optimize cross-entropy at the level of response tokens.
Three matched Llama 3.1 8B seed pairs share candidate membership, input order,
example schedule, base checkpoint and optimization settings. They use the
compact-study frozen 4,494-list artifact. All checkpoints rerank the same LitSearch pools, and
quality is evaluated after the same output repair. Detailed training, inference
and paired-statistical procedures are in Online Resource~1, Section~S15.

\begin{table}[ht]
\centering
\footnotesize
\setlength{\tabcolsep}{4pt}
\renewcommand{\arraystretch}{1.08}
\begin{adjustbox}{max width=\linewidth}
\begin{tabular}{l l c c c}
\hline
Retriever & Tie policy & R@20 & R@5 & nDCG@20 \\
\hline
BM25 & Deterministic & $67.21\pm0.70$ & $54.50\pm0.50$ & $49.89\pm0.30$ \\
     & Epoch-randomized & $67.21\pm0.22$ & $56.17\pm0.25$ & $47.74\pm1.26$ \\
\hline
GRIT & Deterministic & $82.57\pm0.80$ & $69.78\pm0.05$ & $64.79\pm0.16$ \\
     & Epoch-randomized & $80.52\pm0.52$ & $60.98\pm1.86$ & $53.33\pm2.10$ \\
\hline
Qwen3 & Deterministic & $80.94\pm0.73$ & $69.92\pm0.24$ & $64.84\pm0.24$ \\
      & Epoch-randomized & $79.37\pm0.52$ & $62.63\pm2.06$ & $53.26\pm1.93$ \\
\hline
\end{tabular}
\end{adjustbox}
\caption{Matched Llama 3.1 8B tie policies: mean $\pm$ sample standard deviation over three seeds, with scores multiplied by 100. Quality is evaluated after the same repair rule}
\label{tab:tie-sensitivity}
\end{table}

Randomizing ties lowers nDCG@20 with every retriever
(Table~\ref{tab:tie-sensitivity}). Averaged over seeds and retrievers, R@5 and
nDCG@20 decline by 4.81 and 8.40 pp, respectively (Holm-adjusted $p=.0004$
for both); the R@20 difference is not significant ($p_{\mathrm H}=.1404$).
The raw exact-permutation rate falls from 52.67\% to 0.24\%, and the mean number of missing
IDs increases from 1.42 to 32.25. These validity changes qualify the repaired
quality scores. Full intervals and the saved-response output-budget audit are
reported in Online Resource~1, Section~S15.

The experiment establishes sensitivity to target ordering under this transfer
and repair protocol. It does not separate relevance learning from output
validity or transfer from 50/51 training candidates to 100 evaluation candidates.
Epoch-wise randomization still trains against one serialized target; a loss
that accepts alternative within-grade orders remains untested.

\subsection{Comparison with open reranking baselines}
RankZephyr~\cite{Pradeep} and Qwen3-Reranker-8B~\cite{Zhang} receive the same
LitSearch title-abstract inputs, frozen pools and metric code as our models.
Their native interfaces differ: RankZephyr uses RankLLM sliding windows with
context size 4096, window size 20 and stride 10; Qwen3-Reranker-8B uses a local
score-based endpoint. Table~\ref{tab:open-reranker-baselines} compares the systems.

\begin{table}[htbp]
\centering
\footnotesize
\setlength{\tabcolsep}{4pt}
\begin{adjustbox}{max width=\linewidth}
\begin{tabular}{l l c c c}
\hline
Reranker & Retriever & R@20 & R@5 & nDCG@20 \\
\hline
\multicolumn{5}{l}{\textit{Retrieval only}} \\
- & GRIT  & 80.0 & 69.1 & 62.1 \\
- & Qwen3 & 80.5 & 69.0 & 62.9 \\
\hline
\multicolumn{5}{l}{\textit{Reranking on the same frozen candidate pools}} \\
RankZephyr & GRIT  & 81.1 & 69.5 & 64.4 \\
RankZephyr & Qwen3 & 81.5 & 69.4 & 65.1 \\
Qwen3-Reranker-8B & GRIT  & 80.2 & 60.2 & 54.0 \\
Qwen3-Reranker-8B & Qwen3 & 79.2 & 59.7 & 53.0 \\
Gemma 3 27B FT (ours) & GRIT  & 82.8 & 70.9 & 65.0 \\
Gemma 3 27B FT (ours) & Qwen3 & \textbf{84.0} & \textbf{72.5} & \textbf{65.9} \\
\hline
\end{tabular}
\end{adjustbox}
\caption{Local open-reranker comparisons on the same 597 LitSearch queries and frozen top-100 pools. Scores are percentages; bold marks column maxima}
\label{tab:open-reranker-baselines}
\end{table}

RankZephyr reaches approximately 81--82\% R@20 and 64--65\% nDCG@20 with
dense retrieval, close to our stronger configurations. Gemma 3 27B FT retains
higher R@20 and R@5, particularly with Qwen3. Its Qwen3 R@20 advantage over
retrieval alone, RankZephyr and Qwen3-Reranker-8B remains significant after
correction. With GRIT, it remains significant only against retrieval alone.
The nDCG@20 advantage over RankZephyr is not significant with either retriever;
RankZephyr remains statistically competitive on early ordering. The larger
difference from Qwen3-Reranker-8B remains significant. Full paired results are
in Online Resource~1, Table~S8.

\subsection{External GPT-4o reference comparison}
\label{supp:gpt4o}

To position the fine-tuned open-weight models against a strong closed-source
reference, we use the aggregate title-abstract results reported for GPT-4o by
Ajith et al.~\cite{Ajith}. GPT-4o was not rerun. Matching retriever families
therefore provides descriptive context without paired significance tests.

The LitSearch authors report average Broad R@$20$ and average Specific R@$5$,
but not a combined Avg.\ Specific R@20. They do provide results for both specific
subsets. We therefore calculate their query-weighted mean using the weighting
underlying the LitSearch values: $n=231$ inline-specific and $n=211$
author-specific queries. This gives 85.9\% for GPT-4o + GritLM and 74.9\% for
GPT-4o + BM25. We compute our Specific R@20 results in the same way.

\begin{table}[ht]
\centering
\renewcommand{\arraystretch}{1.1}
\setlength{\tabcolsep}{4pt}
\footnotesize
\begin{adjustbox}{max width=\linewidth}
\begin{tabular}{l c c c}
\hline
Configuration & Broad R@20 & Spec.\ R@5 & Spec.\ R@20 \\
\hline
\multicolumn{4}{l}{\textit{Closed-source reference~\cite{Ajith}}} \\
GPT-4o reranking (BM25)        & 59.9 & 68.0 & 74.9 \\
GPT-4o reranking (GritLM)      & 75.3 & 79.2 & 85.9 \\
\hline
\multicolumn{4}{l}{\textit{BM25 first-stage (ours)}} \\
Llama 3.1 8B FT (BM25)         & 53.2 & 57.9 & 70.8 \\
Gemma 3 12B FT (BM25)          & 53.8 & 58.7 & 70.5 \\
Qwen 2.5 14B FT (BM25)         & 48.6 & 57.4 & 70.9 \\
Gemma 3 27B FT (BM25)          & 52.1 & 62.0 & 71.2 \\
\hline
\multicolumn{4}{l}{\textit{GritLM first-stage (ours)}} \\
Llama 3.1 8B FT (GRIT)     & 70.4 & 72.3 & 82.6 \\
Gemma 3 12B FT (GRIT)      & 72.4 & 75.6 & 84.3 \\
Qwen 2.5 14B FT (GRIT)     & 73.3 & 75.5 & 83.3 \\
Gemma 3 27B FT (GRIT)      & 73.0 & 75.8 & \textbf{86.2} \\
\hline
\multicolumn{4}{l}{\textit{Strong open reranker baseline}} \\
RankZephyr (GRIT)          & 72.0 & 74.0 & 84.3 \\
RankZephyr (Qwen3)         & 74.6 & 74.4 & 83.9 \\
\hline
\multicolumn{4}{l}{\textit{Retrieval-only comparison}} \\
Qwen3-Embedding-8B & 69.7 & 74.2 & 84.3 \\
\hline
\multicolumn{4}{l}{\textit{Local end-to-end pipeline with stronger first stage}} \\
Gemma 3 27B FT (Qwen3)     & \textbf{77.6} & 76.9 & \textbf{86.3} \\
\hline
\end{tabular}
\end{adjustbox}
\caption{ Comparison with the closed-source GPT-4o reranker reported by~\cite{Ajith}. “Spec. R@20” is the query-weighted average
over inline-specific and author-specific queries, matching the breakdown used by LitSearch. The BM25 and GritLM rows
support the reported first-stage reranker comparisons; RankZephyr is included as a strong open reranker baseline; the Qwen3 rows
are end-to-end pipeline comparisons with a stronger first-stage retriever. Bold marks values that match or exceed the strongest
GPT-4o configuration (GritLM first-stage).}
\label{tab:gpt4o}

\end{table}

Table~\ref{tab:gpt4o} separates retriever-matched reference positioning from
retrieval-only and end-to-end pipeline results. With BM25, GPT-4o remains ahead of all
fine-tuned open rerankers on the three reported metrics. When we match the
GritLM retriever family, fine-tuned Gemma 3 27B slightly exceeds the reported
GPT-4o Specific R@$20$ ($86.2\%$ vs.\ $85.9\%$), but remains lower on Broad
R@$20$ ($-2.3$ pp) and Specific R@$5$ ($-3.4$ pp). The open reranker is therefore
competitive at the broader R@$20$ cutoff, but it does not outperform GPT-4o
uniformly. We cannot compute confidence intervals or significance tests for this
cross-protocol comparison because per-query GPT-4o outputs are unavailable.

RankZephyr provides a second open reranking reference. On the controlled local
GritLM candidate pool, it obtains $72.0\%$ Broad R@$20$, $74.0\%$ Specific
R@$5$ and $84.3\%$ Specific R@$20$. These results place it close to the stronger
fine-tuned open rerankers, but below our Gemma 3 27B FT result on Specific
R@$20$.

The Qwen3 rows in Table~\ref{tab:gpt4o} compare complete pipelines
(Online Resource~1, Fig.~S2). Fine-tuned Gemma 3 27B with the stronger Qwen3 retriever
achieves $77.6\%$ Broad R@$20$ and $86.3\%$ Specific R@$20$, exceeding the
reported GPT-4o + GritLM pipeline on both R@$20$ metrics. This difference reflects
the combined contribution of a stronger first-stage retriever and the fine-tuned
reranker. It is not a pure comparison of reranker quality against GPT-4o.

\subsection{Compact supervised comparison and inference cost}
\label{sec:compact-results}
GTE-ModernBERT is already competitive before DORIS-MAE adaptation
(Table~\ref{tab:compact-litsearch}). With BM25, both versions exceed every LLM
in Table~\ref{tab:litsearch-main} on all three metrics. Adapted GTE reaches
R@20 of $82.56\%$ with GRIT and $82.98\%$ with Qwen3, although primary Gemma
3 27B FT retains higher R@20 and R@5 point estimates with both dense retrievers.

\begin{table}[ht]
\centering
\footnotesize
\setlength{\tabcolsep}{3pt}
\begin{adjustbox}{max width=\linewidth}
\begin{tabular}{l l c c c}
\hline
GTE condition & Retriever & R@20 & R@5 & nDCG@20 \\
\hline
Zero-shot & BM25 & 71.70 & 62.85 & 58.61 \\
Adapted   & BM25 & 71.54 & 62.12 & 58.55 \\
\hline
Zero-shot & GRIT & 82.28 & 67.29 & 64.36 \\
Adapted   & GRIT & 82.56 & 68.51 & 64.66 \\
\hline
Zero-shot & Qwen3 & 81.78 & 68.12 & 64.26 \\
Adapted   & Qwen3 & 82.98 & 69.10 & 65.05 \\
\hline
\end{tabular}
\end{adjustbox}
\caption{GTE-ModernBERT before and after adaptation on the same 597 LitSearch queries and frozen top-100 pools. Scores are percentages; zero-shot denotes the released ranking-pretrained model}
\label{tab:compact-litsearch}
\end{table}

Adapted GTE exceeds primary Llama FT on R@20 by 5.30, 3.14 and 3.48 pp with
BM25, GRIT and Qwen3, respectively. All remain significant after Holm correction.
The nDCG@20 difference is significant only with BM25. None of GTE's LitSearch
adaptation changes remains significant after correction, including its $+1.20$
pp Qwen3 R@20 gain. Much of its effectiveness is already present in the released
ranking-pretrained checkpoint. We do not report paired significance against
Gemma 3 27B for this experiment.

We measure adapted GTE and locally served Llama 3.1 8B FT on 20 fixed Qwen3
queries, each with 100 candidates truncated to 300 words. Models run sequentially
on an otherwise idle H200, with one warmup and three measured repetitions.
We divide median total time by query count. GTE uses prediction batches of
64 pairs. Llama uses vLLM generation. Thus, the comparison concerns the
implemented inference systems, not architecture alone.

\begin{table}[ht]
\centering
\footnotesize
\setlength{\tabcolsep}{3pt}
\begin{adjustbox}{max width=\linewidth}
\begin{tabular}{l r r}
\hline
Measurement & GTE adapted & Llama FT \\
\hline
Time per query (ms) & 61.6 & 2,525.4 \\
Queries per second & 16.24 & 0.396 \\
Documents per second & 1,623.8 & 39.6 \\
Recorded GPU memory (GiB) & 1.63 & 129.57 \\
\hline
\end{tabular}
\end{adjustbox}
\caption{Inference cost on 20 fixed queries. Time is amortized per query. GTE memory is peak PyTorch reserved memory; Llama memory includes the server allocation and uses different accounting}
\label{tab:compact-compute}
\end{table}

GTE is approximately 41 times faster in this microbenchmark. 
Llama emits a complete permutation without repair for 30\% of the
sampled queries, with a mean of 98.3 parsed indices. Its evaluated rankings use
the completion rule above. These diagnostics are not full-dataset validity rates.

\subsection{CSFCube transfer}
Table~\ref{tab:csfcube} gives the released dev and test aggregates. All three
mid-size models improve nDCG\%@20 on both: Llama gains $+0.023$ and $+0.068$,
Gemma 3 12B gains $+0.034$ and $+0.032$, and Qwen 2.5 14B gains $+0.008$ and
$+0.010$. Gemma 3 27B instead declines by $0.027$ and $0.023$, while R@20
is essentially preserved. These fold-level differences are descriptive.

\begin{table}[ht]
\centering
\renewcommand{\arraystretch}{1.1}
\setlength{\tabcolsep}{4pt}
\footnotesize
\begin{adjustbox}{max width=\linewidth}
\begin{tabular}{l l c c c c}
\hline
\multirow{2}{*}{Model} & \multirow{2}{*}{Metric}
 & \multicolumn{2}{c}{Dev split} & \multicolumn{2}{c}{Test split} \\
 & & base & FT & base & FT \\
\hline
\multirow{2}{*}{Llama 3.1 8B}
  & nDCG\%@20 & 0.355 & \textbf{0.378} & 0.306 & \textbf{0.374} \\
  & R@20      & 0.309 & \textbf{0.393} & 0.250 & \textbf{0.374} \\
\hline
\multirow{2}{*}{Gemma 3 12B}
  & nDCG\%@20 & 0.369 & \textbf{0.403} & 0.347 & \textbf{0.378} \\
  & R@20      & 0.284 & \textbf{0.332} & 0.288 & \textbf{0.330} \\
\hline
\multirow{2}{*}{Qwen 2.5 14B}
  & nDCG\%@20 & 0.372 & \textbf{0.380} & 0.345 & \textbf{0.356} \\
  & R@20      & \textbf{0.311} & 0.301 & 0.270 & \textbf{0.276} \\
\hline
\multirow{2}{*}{Gemma 3 27B (4-bit)}
  & nDCG\%@20 & \textbf{0.468} & 0.441 & \textbf{0.449} & 0.425 \\
  & R@20      & 0.383 & \textbf{0.414} & 0.387 & \textbf{0.389} \\
\hline
\end{tabular}
\end{adjustbox}
\caption{Primary-cohort CSFCube scores under the released aggregation: dev has 24 cases; test averages the 26- and 24-case folds equally. Bold marks the better base/FT value}
\label{tab:csfcube}

\end{table}
\begin{table}[t]
\centering
\scriptsize
\setlength{\tabcolsep}{4pt}
\renewcommand{\arraystretch}{1.0}
\begin{adjustbox}{max width=\linewidth}
\begin{tabular}{l l c c r c c c}
\hline
Model & Metric & Base & FT & $\Delta$ & 95\% CI & $p$ & $p_{\mathrm{H}}$ \\
\hline
Llama 3.1 8B & R@20 & 0.247 & 0.373 & 0.126 & [0.052, 0.215] & .0015 & .0060 \\
Llama 3.1 8B & nDCG\%@20 & 0.304 & 0.374 & 0.070 & [0.020, 0.123] & .0131 & .0524 \\
Gemma 3 12B & R@20 & 0.289 & 0.330 & 0.041 & [$-$0.025, 0.113] & .2434 & .7301 \\
Gemma 3 12B & nDCG\%@20 & 0.346 & 0.378 & 0.031 & [$-$0.008, 0.073] & .1339 & .4017 \\
Qwen 2.5 14B & R@20 & 0.268 & 0.275 & 0.007 & [$-$0.054, 0.065] & .8319 & 1.0000 \\
Qwen 2.5 14B & nDCG\%@20 & 0.344 & 0.355 & 0.010 & [$-$0.029, 0.050] & .6149 & .6629 \\
Gemma 3 27B & R@20 & 0.387 & 0.388 & 0.001 & [$-$0.059, 0.061] & .9840 & 1.0000 \\
Gemma 3 27B & nDCG\%@20 & 0.448 & 0.425 & $-$0.023 & [$-$0.069, 0.023] & .3315 & .6629 \\
\hline
\end{tabular}
\end{adjustbox}
\caption{CSFCube FT-minus-base differences over 50 pooled query-facet cases, with pointwise 95\% intervals clustered by 34 papers. $p_{\mathrm{H}}$ adjusts the four comparisons per metric; case weighting differs from Table~\ref{tab:csfcube}}
\label{tab:csfcube-uncertainty}
\end{table}

Under the case-weighted, paper-clustered analysis in
Table~\ref{tab:csfcube-uncertainty}, only Llama's R@20 gain remains significant
after correction: $+0.126$, with $p_{\mathrm{H}}=.0060$. Its $+0.070$
nDCG\%@20 gain has an interval excluding zero, but $p_{\mathrm{H}}=.0524$.
The other three models' intervals include zero and adjusted values exceed
$.05$. This supports Llama transfer in broader recall across query paradigms.
The nDCG evidence and other model comparisons remain inconclusive.

Table~\ref{tab:csfcube-external-baselines} places our strongest generative FT primary cohort run
against published baselines by Mysore et al.~\cite{Mysore}. Gemma 3 27B FT outperforms fabs\_tfidfcbow200 and
has higher R@20 than UnSimCSE-BERT, with approximately matching nDCG\%@20,
but remains below SPECTER-ID. We interpret this as cross-benchmark transfer,
not a new state of the art. These external artifacts lack matched per-case
outputs, so the comparisons do not receive paired intervals or tests.

\begin{table}[htbp]
\centering
\footnotesize
\setlength{\tabcolsep}{5pt}
\renewcommand{\arraystretch}{1.1}
\begin{adjustbox}{max width=\linewidth}
\begin{tabular}{l c c}
\hline
Method & R@20 & nDCG\%@20 \\
\hline
fabs\_tfidfcbow200 & 0.231 & 0.314 \\
UnSimCSE-BERT      & 0.344 & 0.426 \\
Gemma 3 27B FT (ours) & 0.389 & 0.425 \\
SPECTER-ID         & 0.521 & 0.572 \\
\hline
\end{tabular}
\end{adjustbox}
\caption{Comparison with published CSFCube test baselines from Mysore et al.~\cite{Mysore} (Table 3), aggregated across facets. Published baseline percentages are expressed as decimals}
\label{tab:csfcube-external-baselines}
\end{table}

The compact baseline is stronger on this benchmark
(Table~\ref{tab:compact-csfcube}). Adapted GTE obtains R@20 of $0.6133$ and
nDCG\%@20 of $0.6156$, compared with $0.5439$ and $0.5860$ before adaptation.
Both versions exceed primary Gemma 3 27B FT's $0.389$ and $0.425$.
In the pooled analysis, adaptation increases R@20 by $0.0697$ and nDCG\%@20
by $0.0297$, with $p_{\mathrm{H}}=.0004$ and $.0003$. Adapted GTE also exceeds
Llama FT on both metrics ($p_{\mathrm{H}}<.001$). Its prior ranking training,
full-abstract query representation and separate environment prevent attributing
the advantage to model size or ranking objective alone.

\begin{table}[ht]
\centering
\footnotesize
\setlength{\tabcolsep}{3pt}
\begin{adjustbox}{max width=\linewidth}
\begin{tabular}{l c c c c}
\hline
\multirow{2}{*}{Facet} & \multicolumn{2}{c}{R@20}
 & \multicolumn{2}{c}{nDCG\%@20} \\
 & Zero-shot & Adapted & Zero-shot & Adapted \\
\hline
Background & 0.5489 & 0.6033 & 0.7160 & 0.7387 \\
Method     & 0.5072 & 0.5758 & 0.4455 & 0.4843 \\
Result     & 0.5746 & 0.6591 & 0.6028 & 0.6301 \\
All        & 0.5439 & 0.6133 & 0.5860 & 0.6156 \\
\hline
\end{tabular}
\end{adjustbox}
\caption{GTE-ModernBERT CSFCube test scores before and after adaptation, using the released equal-fold aggregation. Recall treats grades of at least 2 as relevant}
\label{tab:compact-csfcube}
\end{table}

\subsection{Exact-text overlap sensitivity analysis}
\label{sec:document-disjoint}

Table~\ref{tab:leakage-overlap} establishes disjointness at the query level, but it
also reveals exact document reuse across the scientific corpora. We address this
via a LitSearch sensitivity analysis. We remove all 227 exact
normalized title-abstract matches to documents seen during DORIS-MAE prompt
construction. Then, we retrieve each first-stage list to depth 200, remove the
blacklisted documents and refill the list in its original order until it contains
100 candidates outside the exact-text blacklist. We exclude two queries whose gold relevant document is
blacklisted instead of changing their relevance sets. This leaves 595 queries.
With each reconstructed-environment Llama checkpoint served through the local
vLLM endpoint, we then rerank the same filtered BM25,
GRIT and Qwen3 inputs for the base and standard FT models.

\begin{table}[ht]
\centering
\scriptsize
\setlength{\tabcolsep}{2.7pt}
\renewcommand{\arraystretch}{1.08}
\resizebox{\textwidth}{!}{
\begin{adjustbox}{max width=\linewidth}
\begin{tabular}{l c c c c c c c c c}
\hline
& \multicolumn{3}{c}{R@20} & \multicolumn{3}{c}{R@5}
& \multicolumn{3}{c}{nDCG@20} \\
Retriever & Base & FT & $\Delta$
& Base & FT & $\Delta$
& Base & FT & $\Delta$ \\
\hline
BM25  & 53.90 & \textbf{67.29} & 13.39
      & 39.33 & \textbf{55.47} & 16.13
      & 39.05 & \textbf{50.65} & 11.60 \\
GRIT  & 72.42 & \textbf{82.08} & 9.66
      & 52.84 & \textbf{69.87} & 17.03
      & 52.30 & \textbf{64.89} & 12.59 \\
Qwen3 & 74.25 & \textbf{82.03} & 7.79
      & 53.17 & \textbf{69.36} & 16.20
      & 51.44 & \textbf{65.06} & 13.62 \\
\hline
\end{tabular}
\end{adjustbox}
}

\caption{Filtered exact-text overlap LitSearch results for representative reconstructed Llama 3.1 8B. Scores and FT-minus-base differences are percentages over the same 595 queries}
\label{tab:litsearch-document-disjoint}
\end{table}

\begin{samepage}
As Table~\ref{tab:litsearch-document-disjoint} shows, the FT output is higher than the corresponding base output for all three metrics
under each retriever in the saved filtered rerun. This descriptive result supports
transfer beyond exact document reuse for the representative Llama configuration,
but it does not quantify sampling uncertainty or establish document-disjoint
behavior for every model family.
\par
\end{samepage}

\subsection{Robustness and interpretation}
\label{sec:robustness}
Online Resource~1 retains the full control experiments. Their main findings
qualify the preceding results in several ways.

\textbf{Precision and training seeds.} The Gemma 3 27B 16-bit LoRA control
changes LitSearch nDCG@20 relative to 4-bit QLoRA by $+0.92$, $-0.26$ and
$-0.88$ pp for BM25, GRIT and Qwen3. All intervals include zero and none is
significant after correction. The evaluated configurations do not show statistically significant differences in nDCG@20 after correction and the observed differences do not allow us to isolate the contribution of quantization. This control was not run on CSFCube. Three reconstructed
Llama and Gemma 3 12B runs preserve the direction of standard-recipe improvements.
Llama's retriever-specific nDCG@20 standard deviations range from 0.11 to 0.27
points, reinforcing caution around sub-point ablation differences.

\textbf{Inference depth.} With training length fixed, increasing evaluation
depth from 20 to 50 improves mean R@20 by 3.48, 2.98 and 3.98 pp for Llama,
Gemma 3 12B and Gemma 3 27B. Increasing from 50 to 100 gives smaller gains
of 1.04, 1.28 and 0.45 pp. At depths 10 and 20, R@20 is bounded by candidate
membership and includes every supplied candidate. Identical scores there do
not imply identical ranking quality (Online Resource~1, Section~S7).

\section{Discussion and Limitations}
\label{sec:discussion}
The experiments support data-efficient adaptation, with several qualifications.
The learning curve covers one representative model and three seeds. Its early
plateau does not establish a universally minimal budget. The matched controls
support aspect-level supervision, while an independent benefit from distinct
pool augmentation is not established. The historical and reconstructed cohorts
must remain separate and the compact-model comparison does not isolate total
training supervision, architecture or objective.
The parent-grouped learning-rate study provides a separate Llama refit. It does
not validate the historical list-level search or establish parent-disjoint
selection for the other model families.

Broader recall and early ranking quality also differ. Fine-tuning can include
more relevant papers in the top 20 while moving the earliest relevant paper
downward. Examples are retained in Online Resource~1, Table~S9. Larger
base models leave less room for improvement, and Gemma 3 27B has mixed nDCG
behavior that the precision control does not consistently resolve.
These findings motivate stronger top-rank-weighted
losses, harder negatives and tie-aware objectives.

The external GPT-4o comparison is useful for perspective. However, it remains only descriptive: differences in retrieval
and unavailable per-query outputs prevent controlled or significance-tested
claims of superiority.

Finally, a compact ranking-pretrained cross-encoder is a strong practical
alternative. GTE's LitSearch adaptation gains are not statistically resolved,
whereas its CSFCube gains are. Its lower measured inference cost is specific to
the retained microbenchmark and memory figures use different accounting.
These results support evaluating compact supervised systems alongside generative
LLMs rather than treating model scale as necessary for scientific reranking.

\section{Conclusion}
We show that open instruction-tuned LLMs can be adapted into useful scientific
paper rerankers using DORIS-MAE's released relevance labels, without generating
additional teacher rankings. Aspect decomposition, candidate subsampling and
set-unique augmentation produce 4,494 correlated training lists from 100
human-authored parent queries. Matched controls support aspect-level supervision,
but do not establish an independent benefit from distinct-pool augmentation.

For Llama 3.1 8B, most observed improvement appears with 10--25 parents.
Among the tested
budgets, 100 gives the highest mean R@20, while 25 gives the highest R@5 and
nDCG@20. Thus, 100 parents suffice for the observed transfer but are not shown
to be necessary. Transfer evidence is stronger on LitSearch than CSFCube.

The compact GTE baseline shows that competitive scientific reranking does not
require a large generative model. 
Larger-model and GPT-4o comparisons remain
configuration-dependent, and early ranking quality is an important limitation.

Further work should examine top-rank-weighted losses, harder negatives and
multi-seed controls across additional scientific benchmarks. The code will be
made available at \url{https://github.com/Andrei481/hria-scientific-reranking}.

\FloatBarrier
\backmatter
\bmhead{Funding}
This research is partially supported by the project
``Romanian Hub for Artificial Intelligence - HRIA'',
Smart Growth, Digitization and Financial Instruments Program,
2021--2027, MySMIS no.~351416.
\bmhead{Supplementary information}
Online Resource~1 contains detailed methods, robustness checks, complete supporting results and exploratory variants.

\begin{thebibliography}{99}

\bibitem{Ajith}Ajith, A., Xia, M., Chevalier, A., Goyal, T., Chen, D., \& Gao, T. (2024). LitSearch: A retrieval benchmark for scientific literature search. In Proceedings of the 2024 Conference on Empirical Methods in Natural Language Processing (pp. 15068--15083). doi: \url{https://doi.org/10.18653/v1/2024.emnlp-main.840}.

\bibitem{Wang}Wang, J. A., Wang, K., Wang, X., Naidu, P., Bergen, L., \& Paturi, R. (2023). Scientific document retrieval using multi-level aspect-based queries. Advances in Neural Information Processing Systems, 36, 38404--38419. doi: \url{https://doi.org/10.52202/075280-1668}.

\bibitem{Mysore}Mysore, S., O'Gorman, T., McCallum, A., \& Zamani, H. (2021). CSFCube---A test collection of computer science research articles for faceted query by example. arXiv preprint arXiv:2103.12906. doi: \url{https://doi.org/10.48550/arXiv.2103.12906}.

\bibitem{Sun}Sun, W., Yan, L., Ma, X., Wang, S., Ren, P., Chen, Z., et al. (2023). Is ChatGPT good at search? Investigating large language models as re-ranking agents. In Proceedings of the 2023 Conference on Empirical Methods in Natural Language Processing (pp. 14918--14937). doi: \url{https://doi.org/10.18653/v1/2023.emnlp-main.923}.

\bibitem{Karpukhin}Karpukhin, V., Oguz, B., Min, S., Lewis, P., Wu, L., Edunov, S., et al. (2020). Dense passage retrieval for open-domain question answering. In Proceedings of the 2020 Conference on Empirical Methods in Natural Language Processing (EMNLP) (pp. 6769--6781). doi: \url{https://doi.org/10.18653/v1/2020.emnlp-main.550}.

\bibitem{Khattab}Khattab, O., \& Zaharia, M. (2020). ColBERT: Efficient and effective passage search via contextualized late interaction over BERT. In Proceedings of the 43rd International ACM SIGIR Conference on Research and Development in Information Retrieval (pp. 39--48). doi: \url{https://doi.org/10.1145/3397271.3401075}.

\bibitem{Beltagy}Beltagy, I., Lo, K., \& Cohan, A. (2019). SciBERT: A pretrained language model for scientific text. In Proceedings of the 2019 Conference on Empirical Methods in Natural Language Processing and the 9th International Joint Conference on Natural Language Processing (EMNLP-IJCNLP) (pp. 3615--3620). doi: \url{https://doi.org/10.18653/v1/D19-1371}.

\bibitem{Cohan}Cohan, A., Feldman, S., Beltagy, I., Downey, D., \& Weld, D. S. (2020). SPECTER: Document-level representation learning using citation-informed transformers. In Proceedings of the 58th Annual Meeting of the Association for Computational Linguistics (pp. 2270--2282). doi: \url{https://doi.org/10.18653/v1/2020.acl-main.207}.

\bibitem{Thakur}Thakur, N., Reimers, N., Rücklé, A., Srivastava, A., \& Gurevych, I. (2021). BEIR: A heterogenous benchmark for zero-shot evaluation of information retrieval models. arXiv preprint arXiv:2104.08663. doi: \url{https://doi.org/10.48550/arXiv.2104.08663}.

\bibitem{Muennighoff}Muennighoff, N., Tazi, N., Magne, L., \& Reimers, N. (2023). MTEB: Massive text embedding benchmark. In Proceedings of the 17th Conference of the European Chapter of the Association for Computational Linguistics (pp. 2014--2037). doi: \url{https://doi.org/10.18653/v1/2023.eacl-main.148}.

\bibitem{Xia}Xia, F., Liu, T. Y., Wang, J., Zhang, W., \& Li, H. (2008). Listwise approach to learning to rank: Theory and algorithm. In Proceedings of the 25th International Conference on Machine Learning (pp. 1192--1199). doi: \url{https://doi.org/10.1145/1390156.1390306}.

\bibitem{Nogueira}Nogueira, R., Jiang, Z., Pradeep, R., \& Lin, J. (2020). Document ranking with a pretrained sequence-to-sequence model. In Findings of the Association for Computational Linguistics: EMNLP 2020 (pp. 708--718). doi: \url{https://doi.org/10.18653/v1/2020.findings-emnlp.63}.

\bibitem{Pradeep}Pradeep, R., Sharifymoghaddam, S., \& Lin, J. (2023). RankZephyr: Effective and robust zero-shot listwise reranking is a breeze! arXiv preprint arXiv:2312.02724. doi: \url{https://doi.org/10.48550/arXiv.2312.02724}.

\bibitem{Reddy}Gangi Reddy, R., Doo, J., Xu, Y., Sultan, M. A., Swain, D., Sil, A., et al. (2024). FIRST: Faster improved listwise reranking with single token decoding. In Proceedings of the 2024 Conference on Empirical Methods in Natural Language Processing (pp. 8642--8652). doi: \url{https://doi.org/10.18653/v1/2024.emnlp-main.491}.

\bibitem{Zhu}Zhu, Y., Yuan, H., Wang, S., Liu, J., Liu, W., Deng, C., et al. (2026). Large language models for information retrieval: A survey. ACM Transactions on Information Systems, 44(1), Article 12, 1--54. doi: \url{https://doi.org/10.1145/3748304}.

\bibitem{Zhang}Zhang, Y., Li, M., Long, D., Zhang, X., Lin, H., Yang, B., et al. (2025). Qwen3 Embedding: Advancing text embedding and reranking through foundation models. arXiv preprint arXiv:2506.05176. doi: \url{https://doi.org/10.48550/arXiv.2506.05176}.

\bibitem{Hu}Hu, E. J., Shen, Y., Wallis, P., Allen-Zhu, Z., Li, Y., Wang, S., et al. (2022). LoRA: Low-rank adaptation of large language models. In International Conference on Learning Representations (ICLR 2022). arXiv version doi: \url{https://doi.org/10.48550/arXiv.2106.09685}.

\bibitem{Dettmers}Dettmers, T., Pagnoni, A., Holtzman, A., \& Zettlemoyer, L. (2023). QLoRA: Efficient finetuning of quantized LLMs. Advances in Neural Information Processing Systems, 36, 10088--10115. doi: \url{https://doi.org/10.52202/075280-0441}.

\bibitem{Zhuang}Zhuang, H., Qin, Z., Jagerman, R., Hui, K., Ma, J., Lu, J., et al. (2023). RankT5: Fine-tuning T5 for text ranking with ranking losses. In Proceedings of the 46th International ACM SIGIR Conference on Research and Development in Information Retrieval (pp. 2308--2313). doi: \url{https://doi.org/10.1145/3539618.3592047}.

\bibitem{gemma3}Gemma Team, Kamath, A., Ferret, J., Pathak, S., Vieillard, N., Merhej, R., et al. (2025). Gemma 3 technical report. arXiv preprint arXiv:2503.19786. doi: \url{https://doi.org/10.48550/arXiv.2503.19786}.

\bibitem{qwen25}Yang, A., Yang, B., Zhang, B., Hui, B., Zheng, B., Yu, B., et al. (2024). Qwen2.5 technical report. arXiv preprint arXiv:2412.15115. doi: \url{https://doi.org/10.48550/arXiv.2412.15115}.

\bibitem{llama3}Grattafiori, A., Dubey, A., Jauhri, A., Pandey, A., Kadian, A., Al-Dahle, A., et al. (2024). The Llama 3 herd of models. arXiv preprint arXiv:2407.21783. doi: \url{https://doi.org/10.48550/arXiv.2407.21783}.

\bibitem{unsloth}Han, D., Han, M., \& Unsloth Team. (2023). Unsloth [Software]. \url{https://github.com/unslothai/unsloth}.

\bibitem{Kwon}Kwon, W., Li, Z., Zhuang, S., Sheng, Y., Zheng, L., Yu, C. H., et al. (2023). Efficient memory management for large language model serving with PagedAttention. In Proceedings of the 29th Symposium on Operating Systems Principles (pp. 611--626). doi: \url{https://doi.org/10.1145/3600006.3613165}.

\bibitem{gritlm}Muennighoff, N., Su, H., Wang, L., Yang, N., Wei, F., Yu, T., et al. (2025). Generative representational instruction tuning. In International Conference on Learning Representations (ICLR 2025). arXiv version doi: \url{https://doi.org/10.48550/arXiv.2402.09906}.

\bibitem{gte-modernbert}Tongyi Lab, Alibaba Group. (n.d.). GTE-reranker-modernbert-base [Model card]. \url{https://huggingface.co/Alibaba-NLP/gte-reranker-modernbert-base}. Evaluated revision: \texttt{204bdd4}.

\end{thebibliography}

\end{document}
